\section{Специальные функции и бесконечное произведение}

\subsection{Бесконечное произведение. Формула Стирлинга}

\begin{definition}
	Если $c_1,...c_n,...\in\R\backslash\{0\}$, то бесконечным произведением $\prod\limits_{n=1}^\infty c_n$ называют ненулевой предел $\liml_{n\to\infty}P_n, P_n = \prod\limits_{k=1}^n c_k$. Если таковой существует, то бесконечное произведение сходится, в противном случае расходится.
\end{definition}

\begin{theorem} (Необходимое условие сходимости бесконечного произведения)
	Если $\prod\limits_{n=1}^\infty c_n$ сходится, то $\exists \liml_{n \to \infty} c_n = 1$.
\end{theorem}

\begin{proof}
	Коль скоро ряд сходится, это означает, что
	\[
	\exists \liml_{n \to \infty} P_n = P
	\]
	Заметим следующее соотношение:
	\[
	c_n = P_n / P_{n - 1}
	\]
	А отсюда имеем требуемое:
	\[
	\liml_{n \to \infty} c_n = \liml_{n \to \infty} P_n / \liml_{n \to \infty} P_{n - 1} = P / P = 1
	\]
\end{proof}

\begin{theorem}
	(Критерий сходимости бесконечного произведения) Бесконечное произведение $\prod\limits_{n=1}^\infty (1+a_n), a_n > -1$ сходится тогда и только тогда, когда сходится ряд $\sum\limits_{n=1}^\infty \ln{(1 + a_n)}$
\end{theorem}

\begin{proof}
	Получается очевидно логарифмированием префикса произведения и перевода его в префикс суммы (ну и мы пользуемся непрерывностью логарифма).
\end{proof}

\begin{definition}
	Бесконечное произведение сходится абсолютно, если абсолютно сходится соответствующий ему  $\sum\limits_{n=1}^\infty \ln{(1 + a_n)}$ из прошлой теоремы. Несложно видеть, что т.к. в абсолютно сходящемся ряду можно переставлять члены, то то же касается и произведения.
\end{definition}

\begin{theorem}
	(Критерий абсолютной сходимости бесконечного произведения) Бесконечное произведение $\prod\limits_{n=1}^\infty (1+a_n), a_n > -1$ абсолютно сходится тогда и только тогда, когда абсолютно сходится ряд $\sum\limits_{n=1}^\infty a_n$
\end{theorem}

\begin{proof}
	По определению абсолютная сходимость записывается, как абсолютная сходимость ряда $\sum\limits_{n=1}^\infty \ln{(1 + a_n)}$, т.е. сходимость ряда $\sum\limits_{n=1}^\infty |\ln{(1 + a_n)|}$. Заметим, что если необходимое условие не выполняется, то оба объекта из условия расходятся, значит пусть $a_n$ стремятся к нулю, отсюда получаем, что $|\ln{(a_n + 1)}|\sim|a_n|$, что означает эквивалентность сходимости последнего ряда и ряда $\sum\limits_{n=1}^\infty |a_n|$, а это то что нам и нужно.
\end{proof}

\begin{lemma}
	$(\forall n\in\N_0)\emb\sin{x} = (2n + 1)\sin{\frac{x}{2n+1}}\prod\limits_{j=1}^n (1 - \frac{\sin^2{\frac{x}{2n+1}}}{\sin^2{\frac{\pi j}{2n+1}}})$
\end{lemma}

\begin{proof}
	По формуле Муавра $(\cos{x} +i\sin{x})^{2n+1} = \cos{(2n+1)x} +i\sin{(2n+1)x}$ Откуда получается, что для некоторого многочлена $P_n : \deg{P_n} = n$ верно, что $\sin{(2n+1)x} = \sin{x} P_n(\sin^2{x})$. Причём из последней формулы видно, что числа $\sin^2{\frac{j\pi}{2n+1}}, j = 1,...,n$ - его корни, а т.к. все они различны, по основной теореме алгебры это все его корни. Более того, так как
	\[
	\frac{\sin{(2n+1)x}}{\sin{x}} =  P_n(\sin^2{x})
	\], то устремив $x$ к нулю, мы получим, что свободный член $P_n$ это $2n+1$. Теперь зная все корни $P_n$ и его свободный член, мы можем его однозначно выписать:
	\[
	P_n(t) = (2n+1)\prod\limits_{j=1}^n (1-\frac{t}{\sin^2{\frac{j\pi}{2n+1}}})
	\]
	Осталось подставить это в определение $P_n$ :
	\[
	\sin{(2n+1)x} = (2n+1)\sin{x}\prod\limits_{j=1}^n (1-\frac{\sin^2{x}}{\sin^2{\frac{j\pi}{2n+1}}}) 
	\]
	Подстановка $z = (2n+1)x$ теперь даст искомую формулу от $z$ ч.т.д.
\end{proof}

\begin{theorem}
	(Представление синуса бесконечным произведением)
	\[
	\sin{x} = x\prod\limits_{n=1}^\infty (1-\frac{x^2}{n^2\pi^2})
	\]
\end{theorem}

\begin{proof}
	Введём два новых обозначения:
	\[
	U_k^{(n)} := (2n + 1)\sin{\frac{x}{2n+1}}\prod\limits_{j=1}^k (1 - \frac{\sin^2{\frac{x}{2n+1}}}{\sin^2{\frac{\pi j}{2n+1}}})
	\]
	\[
	V_k^{(n)} := \prod\limits_{j=k+1}^n (1 - \frac{\sin^2{\frac{x}{2n+1}}}{\sin^2{\frac{\pi j}{2n+1}}})
	\]
	Зафиксируем $x$. Тогда во-первых, т.к. произведение из условия с некоторого момента имеет близкие к $1$ множители, к нему можно применить критерий абсолютной сходимости, который выполнится, из сходимости ряда обратных квадратов. Во-вторых несложно видеть, что:
	\[
	\exists \liml_{n\to\infty} U_k^{(n)}= x\prod\limits_{j=1}^k (1-\frac{x^2}{j^2\pi^2})=:U_k
	\]
	Т.к. по лемме выше $\sin{x} = U_k^{(n)} V_k^{(n)}$, то из последнего следует также, что
	\[
	\exists \liml_{n\to\infty} V_k^{(n)}=:V_k
	\]
	И соответственно $\sin{x} = U_k V_k$
	Теперь заметим, что если мы покажем, что $\liml_{k\to\infty} V_k = 1$, то, устремив $k$ к бесконечности, мы победим, просто потому что формула из условия, если предел $U_k$ непосредственно. 
	\\
	Итак, покажем, что $\liml_{k\to\infty} V_k = 1$.
	Во-первых легко видеть (просто из взятия производных), что 
	\[
	0<\phi<\frac{\pi}{2}\emb \frac{2}{\pi}\phi<\sin{\phi}<\phi
	\]
	Теперь оценим $V_k^{(n)}$, для достаточно больших $k$, так чтобы аргументы синусов были из ограничений выше:
	\[
	1 > V_k^{(n)} > \prod\limits_{j=k+1}^n (1 - \frac{\frac{x^2}{(2n+1)^2}}{\frac{4}{\pi^2}\frac{\pi^2 j^2}{(2n+1)^2}})= \prod\limits_{j=k+1}^n (1 - \frac{x^2}{4j^2})
	\]
	Опять в силу сходимости ряда обратных квадратов получим, что $\prod\limits_{j=k+1}^\infty (1 - \frac{x^2}{4j^2}) =: W_k$ сходится, и в неравенстве выше можно перейти к пределу: 
	\[
	1\ge V_k\ge W_k
	\]
	А теперь интересный мув. Как мы помним в теории рядов сходимость ряда по сути означает, что суммы его суффиксов с некоторого момента сколь угодно малы. По аналогичным причинам в теории произведений, с некоторого момента суффиксы сколь угодно близки к единице. А $W_k$, как несложно увидеть и есть суффиксы $\prod\limits_{j=1}^\infty (1 - \frac{x^2}{4j^2})$. То есть для достаточно больших $k$:
	\[
	(\forall \eps > 0)\emb 1 \ge V_k \ge W_k > 1 - \eps
	\]
	А это и означает искомое стремление $V_k$ к $1$ ч.т.д.
	
	Здесь надо уточнить, что бесконечное произведение потеряет смысл, если $x = \pi n$, из-за того что она будет 0, но в этом случае формула остаётся верной уже совсем очевидно (если отказаться от смысла большего чем предел произведений префиксов).
\end{proof}

\begin{corollary}
	\[
	\cos{x} = \prod\limits_{n=1}^\infty (1-\frac{4x^2}{(2n-1)^2\pi^2})
	\]
\end{corollary}

\begin{proof}
	Случаи нулевых синуса и косинуса здесь опять руками проверяются, в остальных случаях:
	\[
	\cos{x} =\frac{\sin{2x}}{2\sin{x}} = \frac{2x\prod\limits_{n=1}^\infty (1-\frac{4x^2}{n^2\pi^2})}{2x\prod\limits_{n=1}^\infty (1-\frac{x^2}{n^2\pi^2})} = \cos{x} = \prod\limits_{n=1}^\infty (1-\frac{4x^2}{(2n-1)^2\pi^2})
	\]
	В последнем переходе мы пользуемся абсолютной сходимостью произведений и возможностью переставлять члены обоих внутри и между собой, как и в рядах.
\end{proof}

\begin{corollary}
	(Формула Валлиса)
	\[
	\frac{\pi}{2} = \liml_{n\to\infty}\frac{(2n)!!^2}{(2n-1)!!^2(2n+1)}
	\]
\end{corollary}

\begin{proof}
	Подставим в формулу синуса $x = \frac{\pi}{2}$:
	\[
	\sin{\frac{\pi}{2}} = \frac{\pi}{2}\prod\limits_{n=1}^\infty (1-\frac{\frac{\pi^2}{4}}{n^2\pi^2}) =\frac{\pi}{2} \liml_{n\to\infty} \prod\limits_{m=1}^n \frac{4m^2-1}{4m^2} =\frac{\pi}{2} \liml_{n\to\infty} \frac{(2n-1)!!^2(2n+1)}{(2n)!!^2}
	\]
	Перенос предела в левую часть даёт искомую формулу.
\end{proof}

\begin{lemma}
	(Постоянная Эйлера-Маскерони)
	\[
	\exists\liml_{n\to\infty}(1 + \frac{1}{2}+...+\frac{1}{n} - \ln{n}) = \gamma -\text{постоянная Эйлера-Маскерони}
	\]
\end{lemma}

\begin{proof}
	Обозначим
	\[
	\gamma_n = 1 + \frac{1}{2}+...+\frac{1}{n} - \ln{n}
	\]
	Тогда 
	\[
	\gamma_{n+1} -\gamma_n = \frac{1}{n+1} - \ln{n+1} + \ln{n} = \frac{1}{n+1} - \ln{\frac{n+1}{n}}
	\]
	Покажем, что это меньше $0$. Но это буквально эквивалентно неравенству
	\[
	e < (1+\frac{1}{n})^{n+1}
	\] Последнее верно так как в доказательстве существования числа $e$ мы показывали, что последовательность из правой части стремится к $e$ сверху. Итак, мы показали, что $\gamma_n$ убывают, для сходимости остаётся оценить их снизу. Мы уже знаем, что $\frac{1}{n+1} > \ln{\frac{n+1}{n}}$. Оценим $\gamma_n$ через это:
	\[
	\gamma_n = 1 + \frac{1}{2}+...+\frac{1}{n} - \ln{n}\ge \ln{\frac{2}{1}} + \ln{\frac{3}{2}} +...+\ln{\frac{n+1}{n} - \ln{n}} = \ln{\frac{n+1}{n}} > 0
	\]
	ч.т.д.
\end{proof}

\begin{definition}
	(Гамма-функция)
	\[
	\Gamma(s) = \frac{1}{se^{\gamma^s}}\prod\limits_{n=1}^\infty(1+\frac{s}{n})^{-1}e^{\frac{s}{n}}
	\]
	В этом определении сразу видно только что гамма-функция не определена в целых неположительных.
\end{definition}

\begin{proposition}
	Бесконечное произведение из определения гамма-функции сходится везде кроме целых неположительных (в 0 конечно сходится, но он выкалывается множителем перед произведением, поэтому нам неинтересен).
\end{proposition}

\begin{proof}
	Возьмём логарифм у множителя c достаточно большим номером:
	\[
	|\ln{(1+\frac{s}{n})^{-1}}e^{\frac{s}{n}}| = \left|-\ln{(1+\frac{s}{n})} + \frac{s}{n}\right| = \frac{s^2}{n^2} + o{(\frac{1}{n^2})}
	\]
	Ряд из полученного выражения, как известно, сходится при любом $s$. Значит всё что нужно, это чтобы все множители идущие до тех, у которых взят логарифм были корректны. А это и значит, что $s$ неотрицательно, и знаменатель не обращается в 0.
\end{proof}

\begin{theorem}
	(Формула Эйлера)
	\[
	\Gamma(s)=s^{-1}\prod\limits_{n=1}^\infty(1+\frac{1}{n})^s(1+\frac{s}{n})^{-1}
	\]
	Для всех $s$ кроме целых неположительных.
\end{theorem}

\begin{proof} Из прошлых формул можно получить:
	\[
	\Gamma(s) = \liml_{m\to\infty} s^{-1}e^{-s( 1 + \frac{1}{2}+...+\frac{1}{m} - \ln{m})}\liml_{m\to\infty}\prod\limits_{n=1}^m(1+\frac{s}{n})^{-1}e^{\frac{s}{n}} = \liml_{m\to\infty} s^{-1}e^{-s( 1 + \frac{1}{2}+...+\frac{1}{m} - \ln{m})}\prod\limits_{n=1}^m(1+\frac{s}{n})^{-1}e^{\frac{s}{n}}
	\]
	Все рациональные степени $e$ сокращаются:
	\[
	\Gamma(s)= \liml_{m\to\infty} s^{-1}m^s\prod\limits_{n=1}^m(1+\frac{s}{n})^{-1} = \liml_{m\to\infty} s^{-1}\prod\limits_{n=1}^m\frac{(n+1)^s}{n^s}(1+\frac{s}{n})^{-1} = s^{-1}\prod\limits_{n=1}^\infty(1+\frac{1}{n})^s(1+\frac{s}{n})^{-1}
	\]
\end{proof}

\begin{theorem}
	(Основное свойство гамма-функции/ формула понижения) \\
	Следующие утверждения верны
	\begin{itemize}
		\item $\Gamma(s)s = \Gamma(s + 1), s \neq 0, -1, ...$
		\item $\Gamma(1) = 1$
		
	\end{itemize}
\end{theorem}

\begin{proof}
	Первое:
	\[
	\frac{\Gamma(s+1)}{\Gamma(s)} = \frac{s}{s+1}\prod_{n=1}^\infty\frac{(1+\frac{1}{n})^{s+1}}{(1+\frac{1}{n})^{s}}\frac{1+\frac{s}{n}}{1+\frac{s+1}{n}} = \frac{s}{s+1}\prod\limits_{n=1}^\infty \frac{n+1}{n} \frac{n+s}{n+s+1} = \frac{s}{s+1}\liml_{n\to\infty}\frac{(n+1)(s+1)}{n+s+1} = s
	\]
	Второе:
	\[
	\Gamma(1)=1^{-1}\prod\limits_{n=1}^\infty(1+\frac{1}{n})^1(1+\frac{1}{n})^{-1} = 1
	\]
\end{proof}

\begin{corollary}
	Гамма-функция - обобщение факториала:
	\[
	n! = \Gamma(n + 1)
	\]
\end{corollary}

\begin{proof}
	Очевидно из прошлой теоремы.
\end{proof}

\begin{theorem}
	(Формула дополнения)
	\[
	s\not\in\Z \Rightarrow\ \Gamma(1-s)\Gamma(s) = \frac{\pi}{\sin{\pi s}}
	\]
\end{theorem}

\begin{proof}
	\[
	\Gamma(1-s)\Gamma(s) = (-s)\Gamma(-s)\Gamma(s)=s^{-1}\prod\limits_{n=1}^\infty(1+\frac{1}{n})^{-s}(1+\frac{-s}{n})^{-1}(1+\frac{1}{n})^s(1+\frac{s}{n})^{-1}
	\]
	Сократим лишнее:
	\[
	\Gamma(1-s)\Gamma(s)=s^{-1}\prod\limits_{n=1}^\infty(1+\frac{-s}{n})^{-1}(1+\frac{s}{n})^{-1} = s^{-1}\prod\limits_{n=1}^\infty(1-\frac{s^2}{n^2})^{-1} =  \frac{\pi}{\sin{\pi s}}
	\]
	Последнее верно из формулы синуса.
\end{proof}